% Figure arealaw2d_fill_and_charge.
% Owner: none (note-only figure).
% Used by: gapped-2d-notes AL Sec. 9 (partitions and rounds: one charge).
% Source: arXiv OpenAI2026AreaLaw, sections/08-scanner.tex:95-140 (padded slots, fills,
% fronts (scanner:fronts), charge lists [j-r_0, j+D]) and 200-215 (Lemma
% scanner:histories (2) and (4)).
% Ink-to-index: none (geometry only). Each column is one positive-depth row of one band,
% padded to n = 4 slots (rows 2-5); unsigned depth = column - 1, so column 1 is X = A cap T
% (depth 0); the target edge is straight. Near fills take slots top to bottom, far fills
% bottom to top (the reversed order). Tints are the parts, in the \scanglyph colours of
% shared-spec 4: X red (selected), U grey (neutral), Y violet (collar), V blue (secondary).
% Left panel, after three near fills of depth 3 and two far fills of depth 6: U = depths
% 1, 2 and slots 1-3 of depth 3, so the near front is j_P = 3 (column 4, arrow east); V =
% depth 7 and slots 3-4 of depth 6, so the far front is -j_F = 6 (column 7, arrow west).
% The bracket under the field is the near charge window [j_P-r_0, j_P+D_ch] = depths 2..7
% (columns 3-8), drawn with r_0 = 1, D_ch = 4. The large red dot (centre) is the chosen
% anchor, slot 2 at depth 4, inside the window; the grey contour is its r_0-ball
% {(2,5),(3,4),(3,5),(3,6),(4,5)} (a graph ball, drawn as a cross). The ball meets the
% near side U (at (3,4)) and Y, so the charge moves its four sites still in Y to U (right
% panel; the ball is not redrawn there, its sites are now inside U). The deepest assigned
% site is at depth 5 <= j_P + D_ch/2 = 5, so the history stays good (Lemma (4)), and
% <= j_P + D_ch + r_0 (Lemma (2)). The actual U and Y are no longer depth intervals (Y is
% now in two pieces). Fronts do not move at a charge: the fill schedule is fixed.
% Labels (L4): fronts and window are named once on the left panel, the parts once on the
% right panel (below the field, clear of the window bracket of the left panel).
% Regions (S2): the parts partition the drawn sites; each part is one enclosure (a
% chained rectangle difference, or a site list for the actual U), as in
% al_scan_actual_histories, so the white gaps between tints are between depth rows or
% slots, not missing sites.
% Check opt-out: the joiner \xrightarrow{charge} is a step, not a relation; both panels
% have empty boundary signature and no equality is claimed.
% Boundary signature: empty in both panels.
% Not to scale: r_0 << D_ch << m << L; real depth rows have n slots in a 2-D annulus and
% are not lattice columns; the ball is schematic.
% Label macros: none.
% Fit: C at 5.5mm (\tnshow[5.5mm], about 72 mm wide).
% Related: al_scan_actual_histories (two bands, depth only, no slots; this zooms into
% one band of it); arealaw2d_variable_radii (balls of Prop. trunc). Neither draws a charge.
\begin{tenkzequation}
    \tenkzkernel
    \begin{tenkz}[rows={wire,wire,wire,wire,wire,wire}, cols=8, bonds=none,
            size=s, metrics=compact]
        \tndeclare{species}{centre}{hue=source:red}
        \tn[at=(2,1), skin=dot]{} \tn[at=(2,2), skin=dot]{} \tn[at=(2,3), skin=dot]{} \tn[at=(2,4), skin=dot]{}
        \tn[at=(2,5), skin=dot]{} \tn[at=(2,6), skin=dot]{} \tn[at=(2,7), skin=dot]{} \tn[at=(2,8), skin=dot]{}
        \tn[at=(3,1), skin=dot]{} \tn[at=(3,2), skin=dot]{} \tn[at=(3,3), skin=dot]{} \tn[at=(3,4), skin=dot]{}
        \tn[at=(3,5), skin=dot, species=centre, size=l]{} \tn[at=(3,6), skin=dot]{} \tn[at=(3,7), skin=dot]{} \tn[at=(3,8), skin=dot]{}
        \tn[at=(4,1), skin=dot]{} \tn[at=(4,2), skin=dot]{} \tn[at=(4,3), skin=dot]{} \tn[at=(4,4), skin=dot]{}
        \tn[at=(4,5), skin=dot]{} \tn[at=(4,6), skin=dot]{} \tn[at=(4,7), skin=dot]{} \tn[at=(4,8), skin=dot]{}
        \tn[at=(5,1), skin=dot]{} \tn[at=(5,2), skin=dot]{} \tn[at=(5,3), skin=dot]{} \tn[at=(5,4), skin=dot]{}
        \tn[at=(5,5), skin=dot]{} \tn[at=(5,6), skin=dot]{} \tn[at=(5,7), skin=dot]{} \tn[at=(5,8), skin=dot]{}
        \tnmark[form=enclosure, species=selected, tint]{(2,1) .. (5,1)}{}
        \tnmark[form=enclosure, species=neutral, tint]{(2,2) .. (5,4) - (5,4) .. (5,4)}{}
        \tnmark[form=enclosure, species=collar, tint]{(2,4) .. (5,7) - (2,4) .. (4,4) - (4,7) .. (5,7)}{}
        \tnmark[form=enclosure, species=secondary, tint]{(2,7) .. (5,8) - (2,7) .. (3,7)}{}
        \tnmark[form=enclosure]{{(2,5), (3,4), (3,5), (3,6), (4,5)}}{}
        \tnwire[kind=string, dir=to, name=fcOneNear]{(1,4)}{(1,5)}
        \tnwire[kind=string, dir=to, name=fcOneFar]{(1,7)}{(1,6)}
        \tnmark[form=label, label pos=w]{(1,4)}{$j_P$}
        \tnmark[form=label, label pos=e]{(1,7)}{$-j_F$}
        \tn[at=(6,3), skin=none]{} \tn[at=(6,8), skin=none]{}
        \tnmark[form=bracket, label pos=s]{(6,3) .. (6,8)}{$[j_P-r_0,\,j_P+D_{\rm ch}]$}
    \end{tenkz}
    $\;\xrightarrow{\text{charge}}\;$
    \begin{tenkz}[rows={wire,wire,wire,wire,wire,wire}, cols=8, bonds=none,
            size=s, metrics=compact]
        \tndeclare{species}{centre}{hue=source:red}
        \tn[at=(2,1), skin=dot]{} \tn[at=(2,2), skin=dot]{} \tn[at=(2,3), skin=dot]{} \tn[at=(2,4), skin=dot]{}
        \tn[at=(2,5), skin=dot]{} \tn[at=(2,6), skin=dot]{} \tn[at=(2,7), skin=dot]{} \tn[at=(2,8), skin=dot]{}
        \tn[at=(3,1), skin=dot]{} \tn[at=(3,2), skin=dot]{} \tn[at=(3,3), skin=dot]{} \tn[at=(3,4), skin=dot]{}
        \tn[at=(3,5), skin=dot, species=centre, size=l]{} \tn[at=(3,6), skin=dot]{} \tn[at=(3,7), skin=dot]{} \tn[at=(3,8), skin=dot]{}
        \tn[at=(4,1), skin=dot]{} \tn[at=(4,2), skin=dot]{} \tn[at=(4,3), skin=dot]{} \tn[at=(4,4), skin=dot]{}
        \tn[at=(4,5), skin=dot]{} \tn[at=(4,6), skin=dot]{} \tn[at=(4,7), skin=dot]{} \tn[at=(4,8), skin=dot]{}
        \tn[at=(5,1), skin=dot]{} \tn[at=(5,2), skin=dot]{} \tn[at=(5,3), skin=dot]{} \tn[at=(5,4), skin=dot]{}
        \tn[at=(5,5), skin=dot]{} \tn[at=(5,6), skin=dot]{} \tn[at=(5,7), skin=dot]{} \tn[at=(5,8), skin=dot]{}
        \tnmark[form=enclosure, species=selected, tint]{(2,1) .. (5,1)}{}
        \tnmark[form=enclosure, species=neutral, tint]{{(2,2), (3,2), (4,2), (5,2), (2,3), (3,3), (4,3), (5,3), (2,4), (3,4), (4,4), (2,5), (3,5), (4,5), (3,6)}}{}
        \tnmark[form=enclosure, species=collar, tint]{(2,4) .. (5,7) - (2,4) .. (4,5) - (3,6) .. (3,6) - (4,7) .. (5,7)}{}
        \tnmark[form=enclosure, species=secondary, tint]{(2,7) .. (5,8) - (2,7) .. (3,7)}{}
        \tnwire[kind=string, dir=to, name=fcTwoNear]{(1,4)}{(1,5)}
        \tnwire[kind=string, dir=to, name=fcTwoFar]{(1,7)}{(1,6)}
        \tn[at=(6,1), skin=none]{} \tn[at=(6,3), skin=none]{}
        \tn[at=(6,5), skin=none]{} \tn[at=(6,8), skin=none]{}
        \tnmark[form=label, species=selected, label pos=s]{(6,1)}{$X$}
        \tnmark[form=label, label pos=s]{(6,3)}{$U$}
        \tnmark[form=label, species=collar, label pos=s]{(6,5)}{$Y$}
        \tnmark[form=label, species=secondary, label pos=s]{(6,8)}{$V$}
    \end{tenkz}
\end{tenkzequation}
